\documentclass{ctexart}
\usepackage{template}
\usepackage{esint,extarrows,tikz}
\begin{document}\pagestyle{empty}
\begin{center}\Large
北京大学数学科学学院2024-25高等数学B1期末模拟考试
\end{center}
\begin{problem}[1.(15\songti{分})]
下面的函数极限存在吗?如果存在,请求出其值;如果不存在,请说明理由.
        \begin{enumerate}[label=\tbf{(\arabic*)}]
            \item \textbf{(5\songti{分})}\ $\displaystyle\lim_{x\to0}\left(\dfrac{1}{x}-\dfrac{1}{\tan x}\right)\dfrac{1}{\tan2x}$.
            \item \textbf{(5\songti{分})}\ $\displaystyle\lim_{x\to0}\dfrac{\sqrt[3]{1+x^2+x^4}-1}{\left(\ln(1+x)\right)^2}$.
            \item \textbf{(5\songti{分})}\ $\displaystyle\lim_{(x,y)\to(0,0)}\dfrac{xy}{x^2+(\e^y-1)^2}$.
        \end{enumerate}
\end{problem}
\begin{solution}
\begin{enumerate}[label=\tbf{(\arabic*)}]
        \item 注意到$\tan 2x=\dfrac{2\tan x}{1-\tan^2x}$.于是
            \[\begin{aligned}
                \lim_{x\to0}\left(\dfrac{1}{x}-\dfrac{1}{\tan x}\right)\dfrac{1}{\tan2x}
                &= \lim_{x\to0}\dfrac{(\tan x-x)(1-\tan^2x)}{2x\tan^2 x} \\
                &= \lim_{x\to0}\dfrac{\tan x-x}{2x^3}\cdot\lim_{x\to0}\dfrac{x^2(1-\tan^2x)}{\tan^2x} \\
                &= \lim_{x\to0}\dfrac{\tan^2x}{6x^2}\cdot1^2\\
                &= \dfrac16
            \end{aligned}\]
        \item 注意到$\sqrt[3]{1+x}=1+\dfrac13x+o(x)$,$\ln(1+x)=x+o(x)$.于是
            \[\begin{aligned}
                \lim_{x\to0}\dfrac{\sqrt[3]{1+x^2+x^4}-1}{\left(\ln(1+x)\right)^2}
                &= \lim_{x\to0}\dfrac{1+\dfrac13x^2+\dfrac13x^4+o(x^2)-1}{x^2+o(x^2)} \\
                &= \dfrac13
            \end{aligned}\]
        \item 取$y=kx$.注意到$\e^x=1+x+o(x)$,于是
            \[\lim_{(x,y)\to(0,0)}\dfrac{xy}{x^2+(\e^y-1)^2}=\lim_{(x,y)\to(0,0)}\dfrac{kx^2}{x^2+k^2x^2+o(x^2)}=\dfrac{k}{1+k^2}\]
            于是取不同的$k$对应的路径所得的极限不同,因而原函数极限不存在.
    \end{enumerate}
\end{solution}
\begin{problem}[2.(10\songti{分})]
设欧氏空间$\R^3$中的平面$\Omega$过直线$x+1=y-3=\dfrac{z}2$,且与平面$3x-y-10z=4$垂直.
        \begin{enumerate}[label=\tbf{(\arabic*)}]
            \item \textbf{(5\songti{分})}\ 求$\Omega$的标准方程.
            \item \textbf{(5\songti{分})}\ 已知以原点为球心的球面$S$与$\Omega$相切,求$S$的方程.
        \end{enumerate}
\end{problem}
\begin{solution}
\begin{enumerate}[label=\tbf{(\arabic*)}]
        \item 由题意可知$(3,-1,-10)$与$\Omega$平行.设$\Omega:ax+by+cz+1=0$,于是有
            \[\left\{\begin{array}{l}
                3a-b-10c=0\\
                ax+b(x+4)+c(2x+2)+1\equiv0
            \end{array}\right.\]
            解得$a=\dfrac17,b=-\dfrac27,c=\dfrac1{14}$.于是$\Omega$的方程为$2x-4y+z+14=0$.
        \item 设切点$P(x,y,z)$,$\Omega$的法向量$\vec{u}=(2,-4,1)$与其上一点$M(-7,0,0)$.我们有
            \[\left|\overrightarrow{OP}\right|=\left|\dfrac{\overrightarrow{OM}\cdot\vec{u}}{|\vec{u}|}\right|=\dfrac{14}{\sqrt{21}}=\dfrac{2\sqrt{21}}{3}\]
            于是$S$的半径为$\dfrac{2\sqrt{21}}{3}$,因而其方程为$x^2+y^2+z^2=\dfrac{28}{3}$.
    \end{enumerate}
\end{solution}
\begin{problem}[3.(10\songti{分})]
回答下列问题.
        \begin{enumerate}[label=\tbf{(\arabic*)}]
            \item \textbf{(5\songti{分})}\ 设正整数$n\geqslant3$,$n$元函数$u:\R^n\setminus\{0\}\to\R$满足
                \[u(\li x,n)=\left(\sum_{k=1}^nx_k^2\right)^{\frac{2-n}{2}}\]
                其中$\displaystyle\sum_{k=1}^nx_k^2\neq0$.试求$\displaystyle\sum_{k=1}^{n}u_{x_kx_k}$.
            \item \textbf{(5\songti{分})}\ 设常数$a\in\R$,设$h(x,t)=f(x+at)+g(x-at)$,其中$f:\R\to\R$和$g:\R\to\R$均有连续的二阶导函数.试求$h_{tt}-a^2h_{xx}$.
        \end{enumerate}
\end{problem}
\begin{solution}
\begin{enumerate}[label=\tbf{(\arabic*)},leftmargin=*]
\item 令$S=\sum_{j=1}^nx_j^2>0$.对每个$k$有
\[u_{x_k}=(2-n)x_kS^{-n/2},\qquad
u_{x_kx_k}=(2-n)S^{-n/2-1}(S-nx_k^2)\]
因此$\sum_{k=1}^nu_{x_kx_k}=(2-n)S^{-n/2-1}(nS-nS)=0$.
\item 由链式法则,$h_{tt}=a^2(f''(x+at)+g''(x-at))$,而$h_{xx}=f''(x+at)+g''(x-at)$.
故$h_{tt}-a^2h_{xx}=0$.
\end{enumerate}
\end{solution}
\begin{problem}[4.(10\songti{分})]
求函数$f(x,y)=x^{\sqrt{y}}$在$(1,1)$处的二阶泰勒多项式.
\end{problem}
\begin{solution}
在$(1,1)$处,$f=1$, $f_x=1$, $f_y=0$, $f_{xx}=f_{yy}=0$, $f_{xy}=1/2$.
这里$f_{yy}=f(\ln x)^2/(4y)-f\ln x/(4y^{3/2})$,原解中的第二项符号应为负.
故二阶泰勒多项式为
\[T_2(x,y)=1+(x-1)+\dfrac12(x-1)(y-1)\]
带Peano余项的公式为$f(x,y)=T_2(x,y)+o((x-1)^2+(y-1)^2)$.
\end{solution}
\begin{problem}[5.(10\songti{分})]
设$t\in[0,2\pi]$,$R>0$,$a>0$.求螺旋线$S:\left\{\begin{array}{l}
            x=R\cos t\\
            y=R\sin t\\
            z=at
        \end{array}\right.$在$t=\dfrac\pi4$的切线和法平面的方程.
\end{problem}
\begin{solution}
对应点为$P=(R/\sqrt2,R/\sqrt2,a\pi/4)$,切向量为$\vec\tau=(-R/\sqrt2,R/\sqrt2,a)$.
所以切线的标准方程为
\[\dfrac{x-R/\sqrt2}{-R/\sqrt2}=\dfrac{y-R/\sqrt2}{R/\sqrt2}=\dfrac{z-a\pi/4}{a}\]
法平面方程为
\[-\dfrac R{\sqrt2}\left(x-\dfrac R{\sqrt2}\right)+\dfrac R{\sqrt2}\left(y-\dfrac R{\sqrt2}\right)
+a\left(z-\dfrac{a\pi}4\right)=0\]
\end{solution}
\begin{problem}[6.(10\songti{分})]
设$a,b,c\in\R$.试证明:方程$\e^x=ax^2+bx+c$的互异实根不超过三个.
\end{problem}
\begin{solution}
设$f(x)=ax^2+bx+c-\e^x$.于是$f'(x)=2ax+b-\e^x$,$f''(x)=2a-\e^x$,$f'''(x)=-\e^x$.\\
    不妨设$f(x)=0$有至少四个互异实根,设其为$x_1,\cdots,x_4$,满足$x_1<\cdots<x_4$.\\
    对于任意$x_k,x_{k+1}$(其中$1\leqslant k<4$),有$f(x_k)=f(x_{k+1})$.\\
    根据Rolle中值定理,存在$\xi_k\in(x_k,x_{k+1})$使得$f'(\xi_k)=0$.\\
    于是$f'(x)=0$至少有三个实根$\xi_1,\xi_2,\xi_3$.同理可知$f''(x)=0$至少有两个实根,$f'''(x)=0$至少有一个实根.\\
    而$f'''(x)=-\e^x<0$,即$f'''(x)=0$没有实根,这与假设不符,从而$f(x)=0$至多有三个互异实根.
\end{solution}
\begin{problem}[7.(10\songti{分})]
证明:对于任意给定的$k\in\R$,存在$1$的开邻域$U$和$W$,存在唯一的函数$y=f(x),x\in U,y\in W$满足方程$x^k-3x^2y+3xy^2-y^k=0$.
\end{problem}
\begin{solution}
在$x,y>0$的区域内令$F(x,y)=x^k-3x^2y+3xy^2-y^k$.
有$F(1,1)=0$, $F_y(1,1)=3-k$.
当$k\neq3$时,由隐函数定理,在$(1,1)$附近存在唯一光滑的隐函数.
当$k=3$时,方程恰为$(x-y)^3=0$,唯一解为$y=x$.
两种情形都满足题意;对于非整数$k$,取邻域位于正数范围内,保证实数幂有定义.
\end{solution}
\begin{problem}[8.(15\songti{分})]
求函数$f(x,y,z)=x^3+y^3+z^3-4xyz$在$D=\left\{(x,y,z)\in\R^3:x^2+y^2+z^2\leqslant1\right\}$上的最大值和最小值.
\end{problem}
\begin{solution}
函数在紧球$D$上连续,因此最大、最小值存在.内部驻点满足
\[3x^2=4yz,\qquad3y^2=4xz,\qquad3z^2=4xy\]
若有坐标为零,则全为零;若全不为零,三式相乘给出$27=64$,矛盾.故唯一内部驻点为原点,函数值为$0$.
在单位球面上应用Lagrange乘子法,有$3x^2-4yz=2\lambda x$及其循环式.
两两相减得到$(x-y)(3x+3y+4z-2\lambda)=0$及循环式.
若三个坐标互不相同,三个括号同时为零,相减迫使坐标相等,矛盾.所以至少两个坐标相等.
若$x=y=z$,则各坐标为$\pm1/\sqrt3$,函数值为$\mp1/(3\sqrt3)$.
否则,由对称性不妨令$x=y\neq z$,则$2\lambda=7x+3z$,并且$x(4x+7z)=0$.
若$x=0$,则点为$(0,0,\pm1)$,函数值为$\pm1$.
若$4x+7z=0$,由球面约束得到
\[(x,y,z)=\pm\dfrac{(7,7,-4)}{\sqrt{114}}\]
及其坐标排列,函数值为$\pm37/(3\sqrt{114})$.
比较所有候选值,得到
\[\max_D f=\dfrac{37}{3\sqrt{114}},\qquad\min_D f=-\dfrac{37}{3\sqrt{114}}\]
最大值在$(7,7,-4)/\sqrt{114}$的三个坐标排列处取得,最小值在它们的相反点处取得.
\end{solution}
\begin{problem}[9.(10\songti{分})]
设函数$f$在$[0,1]$二阶可微,且$f(0)=f(1)=0$,$\displaystyle\min_{x\in[0,1]}f(x)=-1$.试证明:存在$\epsilon\in(0,1)$使得$f''(\epsilon)\geqslant8$.
\end{problem}
\begin{solution}
设$f(x)$在$x=x_0$处取到最小值.将$f(x)$在$x=x_0$处做泰勒展开有
    \[f(x)=f(x_0)+f'(x_0)(x-x_0)+\dfrac{f''(\xi)(x-x_0)^2}{2},x\gtrless\xi\gtrless x_0\]
    由于$f(x)$在$[0,1]$二阶可微,因而最小值点必然满足$f'(x_0)=0$.又$f(x_0)=-1$,于是有
    \[f(x)=\dfrac{f''(\xi)(x-x_0)^2}{2}-1,x\gtrless\xi\gtrless x_0\]
    分别取$x=0,1$有
    \[\left\{\begin{array}{l}
        1=\dfrac{f''(\xi_1)x_0^2}{2}\\
        1=\dfrac{f''(\xi_2)(1-x_0)^2}{2}
    \end{array}\right.\]
    由于$\min\left\{x_0,1-x_0\right\}\leqslant\dfrac12$,于是
    \[\max\left\{f''(\xi_1),f''(\xi_2)\right\}=\max\left\{\dfrac{2}{x_0^2},\dfrac{2}{(1-x_0)^2}\right\}\geqslant\dfrac{2}{\left(\frac12\right)^2}=8\]
\end{solution}
\end{document}
